\begin{abstract}
Financial risk assessment plays a vital role in accelerating the global energy transition and achieving the 1.5$^\circ$C climate target. This study develops a three-stage machine learning framework to analyze the impact of fossil fuel price volatility on the weighted average cost of capital (WACC) for energy transition projects. First, a hybrid Long Short-Term Memory-Gaussian Process Regression (LSTM-GPR) model forecasts fossil fuel prices, followed by an LSTM-based inflation predictor. In the final stage, an AutoML (AutoGluon) regression model evaluates how variations in fossil fuel prices and inflation influence WACC across nine energy sectors and multiple global regions.

The results show regional differences in model fit: average $R^2$ is 0.670 for Europe, 0.650 for China, and 0.547 for the Middle East, with stronger fit for solar and wind projects. These fit statistics describe the joint fuel-price and inflation specifications and do not measure the magnitude of price-shock responses. Additional holdout, rolling-origin, and transmission-diagnostic checks are used to interpret these findings as scenario comparisons and mechanism evidence rather than as formal causal identification. The approach integrates AI-driven forecasting with climate finance modeling, providing an interpretable and data-driven tool to support policy design, green investment strategies, and climate-resilient energy planning.
\end{abstract}

\begin{IEEEkeywords}
Weighted average cost of capital, fossil fuel price volatility, climate finance, energy transition, LSTM-GPR, AutoGluon.
\end{IEEEkeywords}
